% !TeX root = ../main.tex
\begin{figure}[H]
\centering
\begin{minipage}[t]{.48\linewidth}
\raggedright\textbf{Source (annotations inferred)}
\begin{lstlisting}[xleftmargin=0pt]
let step = fun path => fun t' =>
  match path with
    | [] => []
    | current :: _ =>
        let t = fst current in
        let y = snd current in
        let y' = gaussian[E](y, t' - t) in
        (t', y') :: path
in
let times =
  1 :: 2 :: 4 :: 5.5 :: 6 :: 6.25 :: []
in
reduce step ((0, 0) :: []) times
\end{lstlisting}
\end{minipage}\hfill
\begin{minipage}[t]{.48\linewidth}
\raggedright\textbf{After replacement (simplified)}
\begin{lstlisting}[xleftmargin=0pt]
let step = fun path => fun t' =>
  match path with
    | [] => []
    | current :: _ =>
        let y = snd current in
        (t', y) :: path
in
let times =
  1 :: 2 :: 4 :: 5.5 :: 6 :: 6.25 :: []
in
reduce step ((0, 0) :: []) times
\end{lstlisting}
\end{minipage}
\par\smallskip
The source returns a random path; after replacement it returns
$[(6.25,0),(6,0),(5.5,0),(4,0),(2,0),(1,0),(0,0)]$.
\caption{A Gaussian random walk from~\cite{vakar2019domain}.
Both programs return the complete list of sampled time--position pairs.
\texttt{reduce} takes a parameterized random
operation, an initial value, and a stream.}
\label{fig:intro-gaussian-walk}
\Description{The source passes a probabilistic step function to the recursive
reduce function. At each requested time, the step draws a Gaussian position
whose mean is the previous position and whose variance is the elapsed time, and
prepends the new time--position pair to the path. Replacing every draw by its
mean leaves the position fixed at zero.}
\end{figure}
